\chapter{AST-Guided Sorrifier}
\label{sec:Sorrifier}

\section{Motivation}

Lean failures are often local. A generated proof may contain a valid prefix and one broken command, but normal binary verification rejects the whole script. The Sorrifier recovers partial signal by replacing localized failing regions with \texttt{sorry}.

We distinguish two notions:
\begin{itemize}
    \item \textbf{Placeholder-compilable script:} Lean accepts the script only because unresolved or invalid regions were replaced with \texttt{sorry}.
    \item \textbf{Fully verified proof:} Lean accepts the script with no \texttt{sorry} placeholder.
\end{itemize}
The Sorrifier is used only for reward estimation and partial-credit extraction. A theorem is considered solved only when Lean accepts the proof without any \texttt{sorry}.

\section{Failure Taxonomy}

Generated Lean failures are categorized conceptually as follows:
\begin{itemize}
    \item \textbf{Parser errors:} malformed syntax, such as missing brackets, invalid tactic syntax, or broken proof blocks.
    \item \textbf{Elaboration errors:} unresolved names, missing implicit arguments, or expressions that Lean cannot elaborate.
    \item \textbf{Type errors:} syntactically valid terms whose inferred types do not match the expected goal.
    \item \textbf{Unsolved goals:} valid proof fragments that execute but leave remaining subgoals.
\end{itemize}

This taxonomy is used to explain the main failure modes observed in generated Lean code. The current
implementation does not rely on a complete classifier that formally separates parser, elaboration, and
type-checking failures. Instead, candidate actions are executed through a persistent Lean worker, while
an AST daemon provides syntax spans when a parse tree is available. The worker returns a structured
result containing diagnostics and proof-state metadata. The main returned fields are
\texttt{errors}, \texttt{warnings}, \texttt{infos}, and \texttt{sorries}. Each diagnostic contains a severity level,
a textual message, and a source position:
\begin{verbatim}
{
  "pos": {"line": ..., "column": ...},
  "endPos": {"line": ..., "column": ...},
  "severity": "error" | "warning" | "info",
  "data": "..."
}
\end{verbatim}

Operationally, the implementation performs only a coarse separation between unresolved-goal cases and
fatal failures. Diagnostics indicating remaining goals are handled together with the \texttt{sorries} field,
which provides the proof states associated with placeholder locations. Other diagnostic failures are passed
to the Sorrifier as recoverable fatal errors. Worker crashes, system failures, and timeouts are not
repair targets; they are logged as failed actions.

The recovery pipeline is therefore diagnostic-driven rather than taxonomy-driven. It starts from the
first Lean error location, or from the first unsolved-goal location when the candidate elaborates but does
not close the expected target. Given this source position, the Sorrifier attempts to identify the
affected region and replace it with a placeholder. If an AST can be constructed, the system selects the
smallest enclosing tactic, tactic sequence, expression, or declaration block that covers the reported
source position. If the AST is only partial or cannot be constructed due to severe syntax failure, the
system falls back to a coarser line-, indentation-, or block-based recovery strategy. This fallback
searches for nearby structural boundaries such as \texttt{have}, bullet blocks, or \texttt{:= by} proof bodies.

Different patch units are used depending on how precisely the failure can be localized. Local errors may
be patched at the level of a single tactic line or AST node. Repeated failures or severe malformed blocks
may trigger a larger structural patch that removes or replaces an enclosing command block. If the failure
occurs in a declaration body and local recovery is unreliable, the system may replace the body with a
placeholder proof:
\begin{verbatim}
:= by
  sorry
\end{verbatim}

Unsolved goals are handled differently from fatal failures. For skeleton actions, unresolved goals reported
through \texttt{sorries} are converted into child OR nodes in the AND/OR graph. For local proof
actions, remaining goals mean that the action did not discharge the current node. Such actions may still
receive syntactic recovery reward, but they are not allowed to expand the graph. A more fine-grained
classifier for parser, elaboration, and type errors is left as a possible implementation refinement.

\section{AST-based Error Localization}

When Lean reports an error source span for a generated script, the Sorrifier maps that span to the narrowest enclosing syntax node in the script's AST. The preferred repair units are a tactic, a tactic sequence, an expression, or a declaration proof body. The system patches the smallest viable unit first, preserving the largest possible amount of surrounding generated code.

If the error belongs to a declaration body, the body can be replaced by a placeholder proof:
\begin{verbatim}
:= by
  sorry
\end{verbatim}
If the error belongs to a local tactic or expression, only that local region is replaced where possible.

\paragraph{Iterative Patching Algorithm}

The recovery loop operates as follows:
\begin{enumerate}
    \item Submit the generated script to Lean for elaboration and verification.
    \item If the proof is fully verified without placeholders, return full success.
    \item If Lean reports a syntax error, elaboration error, type error, or unsolved-goal span, locate the smallest enclosing \texttt{Syntax} node or structural block for the diagnostic coordinates.
    \item Replace the failing node with \texttt{sorry} or a placeholder proof block.
    \item Re-submit the patched script to Lean.
    \item Stop when the script becomes placeholder-compilable or the patch budget is exhausted.
\end{enumerate}

To guarantee termination and avoid cyclic oscillation, the implementation stores hashes of intermediate patched scripts. If a repeated broken state is detected, the algorithm escalates the patch boundary to a larger enclosing structural block, such as an entire \texttt{have} statement or tactic block. If no reliable local repair remains, the action is marked as failed for search while retaining diagnostics for offline analysis.

The iterative nature of the algorithm matters because Lean errors often cascade. A single type mismatch
or missing argument can make later lines elaborate in the wrong local context, and an unknown identifier
can cause subsequent tactics to report misleading errors. Patching only the first reported error in one
pass is therefore not always enough. GammaZero repeatedly submits the patched candidate to Lean so
that each next diagnostic is computed in the updated proof state. This makes the recovery process
conservative: each structural replacement is justified by an active Lean response rather than by a purely
textual guess about which lines look suspicious.

Crucially, the repair loop is not allowed to become a proof search procedure of its own. It does not
synthesize new lemmas, search for alternative theorem bindings, or replace broken lines with
semantically different tactics. Its role is to preserve the largest checkable prefix and surrounding
structure by isolating failures behind placeholders. This keeps the reward interpretation simple. A high
syntactic score means that much of the model's original code remained Lean-checkable after local
placeholder insertion; it does not mean that the repair loop invented a new proof path.

\section{Syntactic Survival Signal}

After recovery, the system estimates how much of the original generated proof body survives in the
patched, placeholder-compilable script. Although AST information may be used to localize failing
regions during recovery, the current reward implementation does not count AST nodes directly.
Instead, it uses a line-based syntactic score over cleaned proof lines.

The proof body is normalized by removing blank lines and comment-only lines. The remaining lines
from the original generation are compared with the lines in the patched script using sequence matching.
Lines that remain unchanged and in the same relative order are counted as surviving. Lines removed
by patching, deletion, or orphan-code cleanup are treated as non-surviving.

This produces a syntactic signal: generations that require only small localized patches
receive higher syntactic reward, while generations whose proof bodies must be mostly replaced by \texttt{sorry}
receive lower syntactic reward. The exact formula, including the penalty for unused or ineffective code diagnostics,
is defined in Section~\ref{sec:reward}.

The line-based score is deliberately coarse. Counting AST nodes would make the reward depend on parser details rather than on the proof text produced by the model. 
Counting surviving lines makes the reward easier to inspect: the recorded trajectory shows exactly which generated lines were preserved and which lines were replaced by placeholders. 
A short proof whose only line fails should receive little syntactic credit; a longer skeleton that contains several valid declarations before a
localized failure should receive more. This is not a proof-quality metric, but it is useful feedback for
future optimization because it distinguishes completely unusable generations from generations that are
nearly well formed.

There are also cases where recovery is intentionally disabled. Timeouts, worker crashes, and system-level
failures do not identify a trustworthy local source span, so replacing text with \texttt{sorry} would create a
misleading syntactic reward signal. Similarly, action-policy violations such as using \texttt{admit} in a skeleton are
not treated as ordinary Lean failures. They indicate that the model did not follow the action requirement and
are recorded as rejected candidates.

A limitation of the syntactic score is that it does not measure semantic importance. A single surviving
line might introduce a crucial subgoal, while several surviving lines might only set up unused
local names. This is why \(r_{env}\) is later paired with the dependency reward \(r_{dep}\), rather than used alone. The
syntactic reward answers a narrow question: how much generated Lean structure survived recovery?
The downstream dependency reward \(r_{dep}\) answers a later question: did that structure matter after the proof was stitched
and elaborated? Keeping these rewards separate makes the trajectory more interpretable.

Another limitation is that recovery may overestimate low-quality generations. If a long invalid proof contains a few surviving fragments, the line-based score can still assign nonzero reward even when the overall proof direction is poor. 
A model that repeatedly emits long invalid scripts may still obtain nonzero syntactic rewards if some prefixes
parse. For this reason, failed actions do not solve states and repaired skeletons receive a mild priority
penalty. The search can learn from repaired code, but it still prefers clean actions when other signals are
similar.


Concrete Sorrifier examples and line-based reward calculations are moved to Appendix~\ref{app:sorrifier-examples}.
